\section{Subset Sums II (Unique Subsets with Duplicates Present)}
\label{sec:rec-subset-sum-2}

\problemstatement{Given an array $\textit{arr}$ of $n$ integers that
\textbf{may contain duplicate values}, return every distinct subset
(the power set), with \textbf{no duplicate subsets} in the output, even
though duplicate elements exist in the input.}

\begin{patternbox}
The keyword \emph{``may contain duplicates,''} combined with
\emph{``no duplicate subsets in the output,''} is the trigger for this
chapter's other major backtracking template, distinct from pick/not-pick:
\textbf{sort first, then iterate through choices at each recursive level,
explicitly skipping any choice equal to the immediately previous choice
already tried at the \emph{same} level.} This ``sort, then skip
same-level duplicates'' technique is reused in every later problem that
needs to handle duplicate input values
(Sections~\ref{sec:rec-combination-sum-2}--\ref{sec:rec-combination-sum-3}).
\end{patternbox}

\subsection*{Understanding the problem carefully}

If the input has duplicate values, the plain pick/not-pick tree from
Section~\ref{sec:rec-power-set} would generate the \emph{same} subset
multiple times -- for example, with $\textit{arr}=[1,2,2]$, picking
``the first $2$'' and picking ``the second $2$'' (while skipping the
other) produce two different leaves in the recursion tree, but the
exact same subset $\{1,2\}$. To avoid this, we change the recursion's
shape: instead of a binary pick/not-pick at every fixed index, we
\textbf{sort} the array first (so equal values become adjacent), and at
each recursive level, we loop through the remaining choices, but skip
any choice that is equal to the choice immediately before it \emph{at
this same level} (not across different levels, where repeats are
exactly what we want, to build subsets containing multiple copies of a
repeated value).

\begin{ideabox}[Sort First, Then Skip Same-Level Duplicates]
Sort $\textit{arr}$. $\textsc{Solve}(\textit{start}, \textit{current})$:
record a copy of $\textit{current}$ as one valid subset immediately
(every prefix explored is itself a valid, distinct subset, not just the
leaves). For $i$ from $\textit{start}$ to $n-1$: if $i > \textit{start}$
and $\textit{arr}[i] = \textit{arr}[i-1]$, skip this $i$ (this exact
value has already been tried as the \emph{first} new element added at
this level, so trying it again here would only reproduce a subset
already generated). Otherwise: append $\textit{arr}[i]$ to
$\textit{current}$; recurse $\textsc{Solve}(i+1, \textit{current})$;
remove $\textit{arr}[i]$ from $\textit{current}$ (undo).
\end{ideabox}

\subsection*{Why skipping only at the same level (not across levels) is correct}

The key distinction: $\textit{arr}[i] = \textit{arr}[i-1]$ being skipped
only applies when both are being considered as candidates \emph{to fill
the same ``next new element'' slot} within the same call to
$\textsc{Solve}$ -- i.e.\ within the same \texttt{for} loop, at the same
recursion depth. A repeated value is never forbidden from appearing
\emph{multiple times within one subset} (e.g.\ $\{2,2\}$ from
$\textit{arr}=[1,2,2]$ is a perfectly valid, distinct subset) -- what is
forbidden is generating the \emph{same subset twice} by choosing ``the
first $2$ at this position'' in one branch and ``the second $2$ at this
same position'' in a sibling branch, when both choices would lead to
identical subsequent subsets. Since the array is sorted, any two equal
values are guaranteed to be adjacent in the loop, so checking
$\textit{arr}[i] = \textit{arr}[i-1]$ exactly identifies this situation.

\subsection*{Algorithm}

\begin{enumerate}
  \item Sort $\textit{arr}$.
  \item $\textsc{Solve}(\textit{start}, \textit{current})$:
  \begin{enumerate}
    \item Append a copy of $\textit{current}$ to the result.
    \item For $i$ from $\textit{start}$ to $n-1$:
    \begin{itemize}
      \item If $i>\textit{start}$ and $\textit{arr}[i] =
            \textit{arr}[i-1]$: skip (continue to next $i$).
      \item Append $\textit{arr}[i]$ to $\textit{current}$; recurse
            $\textsc{Solve}(i+1, \textit{current})$; remove
            $\textit{arr}[i]$ (undo).
    \end{itemize}
  \end{enumerate}
  \item Call $\textsc{Solve}(0, [\,])$.
\end{enumerate}

\subsection*{Dry run}

\dryrun{Take $\textit{arr} = [1,2,2]$ (already sorted).}

\begin{itemize}
  \item $\textsc{Solve}(0,[])$: record $[]$. Loop $i=0$ ($1$): append.
        $\textsc{Solve}(1,[1])$. Undo.
  \begin{itemize}
    \item $\textsc{Solve}(1,[1])$: record $[1]$. Loop $i=1$ ($2$):
          append. $\textsc{Solve}(2,[1,2])$. Undo. $i=2$ ($2$): $i>
          \textit{start}=1$ and $\textit{arr}[2]=\textit{arr}[1]=2$:
          skip.
    \begin{itemize}
      \item $\textsc{Solve}(2,[1,2])$: record $[1,2]$. Loop $i=2$
            ($2$): append. $\textsc{Solve}(3,[1,2,2])$: record
            $[1,2,2]$. Undo.
    \end{itemize}
  \end{itemize}
  \item Loop continues at top level: $i=1$ ($2$): $i>\textit{start}=0$,
        but $\textit{arr}[1] \ne \textit{arr}[0]$ ($2\ne1$), no skip.
        Append. $\textsc{Solve}(2,[2])$. Undo.
  \begin{itemize}
    \item $\textsc{Solve}(2,[2])$: record $[2]$. Loop $i=2$ ($2$):
          append. $\textsc{Solve}(3,[2,2])$: record $[2,2]$. Undo.
  \end{itemize}
  \item $i=2$ at top level: $i>\textit{start}=0$ and
        $\textit{arr}[2]=\textit{arr}[1]=2$: skip.
\end{itemize}

Recorded subsets: $[], [1], [1,2], [1,2,2], [2], [2,2]$ -- exactly the
$6$ distinct subsets of $\{1,2,2\}$ (note there are $6$, not $2^3=8$,
because $\{2\}$ and $\{2,2\}$ each appear in only one form despite the
duplicate $2$).

\subsection*{C++ implementation}

\begin{lstlisting}
#include <bits/stdc++.h>
using namespace std;

void solve(int start, vector<int>& arr, vector<int>& current,
           vector<vector<int>>& result) {
    result.push_back(current);

    for (int i = start; i < (int)arr.size(); i++) {
        if (i > start && arr[i] == arr[i - 1]) continue; // skip duplicate

        current.push_back(arr[i]);
        solve(i + 1, arr, current, result);
        current.pop_back();
    }
}

vector<vector<int>> subsetsWithDup(vector<int>& arr) {
    sort(arr.begin(), arr.end());
    vector<vector<int>> result;
    vector<int> current;
    solve(0, arr, current, result);
    return result;
}

int main() {
    vector<int> arr = {1, 2, 2};
    for (auto& subset : subsetsWithDup(arr)) {
        cout << "[ ";
        for (int x : subset) cout << x << " ";
        cout << "] ";
    }
    cout << endl;
    return 0;
}
\end{lstlisting}

\begin{complexitybox}
\textbf{Time Complexity.} Sorting costs $O(n \log n)$. The recursion
visits each distinct subset once, and building/copying each costs up to
$O(n)$, giving an overall time complexity bounded by
\[
  O(n \cdot 2^n)
\]
in the worst case (no duplicates, so all $2^n$ subsets are distinct).
\textbf{Space Complexity.} $O(n)$ for the recursion depth and
$\textit{current}$ list.
\end{complexitybox}

\begin{takeawaybox}
``Sort first, then skip a choice equal to the immediately previous
choice tried at the same recursion level'' is the standard fix for
duplicate-avoidance in backtracking, and it generalizes far beyond
subsets: the same technique, applied to the choice loop inside
Combination Sum II (Section~\ref{sec:rec-combination-sum-2}) and
Combination Sum III (Section~\ref{sec:rec-combination-sum-3}), prevents
duplicate combinations there too.
\end{takeawaybox}
